% All four files (main.tex, neurips_2026.sty, neurips_2026.tex, checklist.tex)
% must be in the SAME folder (as on Overleaf).  LaTeX will find the .sty automatically.
\documentclass{article}
\usepackage{neurips_2026}   % change to [main,final] for camera-ready

\usepackage[utf8]{inputenc}
\usepackage[T1]{fontenc}
\usepackage{hyperref}
\usepackage{url}
\usepackage{booktabs}
\usepackage{amsfonts}
\usepackage{nicefrac}
\usepackage{microtype}
\usepackage{xcolor}
\usepackage{graphicx}

\title{BrailleLens: Real-Time Braille Reading Assistance via Fingertip-Guided
Cell Detection and Sinhala Transliteration on Mobile Devices}

\author{%
  Author 1 \\
  Department of Computer Science and Engineering\\
  University of Moratuwa\\
  \texttt{author1@cse.mrt.ac.lk}
  \And
  Author 2 \\
  Department of Computer Science and Engineering\\
  University of Moratuwa\\
  \texttt{author2@cse.mrt.ac.lk}
  \And
  Author 3 \\
  Department of Computer Science and Engineering\\
  University of Moratuwa\\
  \texttt{author3@cse.mrt.ac.lk}
  \And
  Author 4 \\
  Department of Computer Science and Engineering\\
  University of Moratuwa\\
  \texttt{author4@cse.mrt.ac.lk}
}

\begin{document}

\maketitle

\begin{abstract}
We present \textbf{BrailleLens}, a mobile-first system that allows visually
impaired readers to learn Braille interactively without sighted assistance.
A user places a finger on an embossed Braille page in front of a smartphone
camera; the system detects the fingertip in real time, maps its position to
the nearest Braille cell, and speaks or displays the corresponding Sinhala
character. The back-end combines a fine-tuned YOLO v2.6 nano fingertip
detector, an ORB+RANSAC homography registration stage, and a 26-class
convolutional neural network (CNN) trained on fused DSBI and Angelina data
(157,559 annotated cells). On a held-out test set the cell classifier
reaches 94.7\% top-1 accuracy; the mobile-quantised fingertip detector runs
at $\geq$15\,FPS on a mid-range Android device. A Flutter application
delivers the full pipeline entirely on-device, requiring no cloud
connectivity.
\end{abstract}

% ======================================================================
\section{Introduction}
% ======================================================================

Braille literacy is a critical pathway to education for the estimated 43
million blind people worldwide~\cite{who2023}, yet access to sighted Braille
tutors is severely limited in low-resource settings such as Sri Lanka.
Existing digital aids either require expensive refreshable Braille displays,
depend on internet connectivity, or do not support Sinhala—the primary
written language of Sri Lanka.

We address this gap with \textbf{BrailleLens}, a system with two
complementary interaction modes: (1)~\emph{auto-scan mode}, in which the
app passively recognises a full page and reads it aloud, and (2)~\emph{fingertip
learning mode}, in which a child traces individual cells while the system
narrates each character in real time.  The fingertip mode is the technical
centrepiece: detecting a fingertip robustly under uncontrolled lighting,
mapping it to the correct Braille cell despite camera shake, and doing so
at interactive frame rates entirely on-device.

Our key contributions are:
\begin{itemize}
  \item A multi-source Braille cell dataset (DSBI + Angelina + Gold,
    157,559 cells) with a reproducible preprocessing pipeline that performs
    \emph{box harmonisation} across annotation conventions.
  \item A domain-adapted YOLO v2.6 nano fingertip detector fine-tuned on
    Braille-reading video, exported to UINT8 ONNX for mobile deployment.
  \item An ORB+RANSAC homography registration approach that tolerates
    moderate page drift without fiducial markers.
  \item A Flutter application delivering the complete pipeline on-device
    for both Android and iOS.
\end{itemize}

% ======================================================================
\section{Related Work}
% ======================================================================

\paragraph{Braille recognition.} Prior optical Braille recognition systems
largely rely on controlled scanning hardware~\cite{antonacopoulos1998}.
Recent deep-learning approaches~\cite{ovodov2021} improve robustness but
target printed Braille from flatbed scans and do not address Sinhala.
The Angelina reader~\cite{angel2021} processes smartphone photos but
requires a back-end server.

\paragraph{Fingertip tracking.} MediaPipe Hands~\cite{zhang2020} provides
reliable skeleton-level detection but struggles with partial hand occlusion
by the page surface.  Domain-specific YOLO detectors have shown superior
recall in constrained interaction scenarios~\cite{redmon2016}.

\paragraph{Mobile Braille tutors.} Commercial apps (e.g.\ Read2Me, KNFB
Reader) focus on OCR of print text; none offer interactive Braille cell
identification with Sinhala output.

% ======================================================================
\section{Methodology}
% ======================================================================

\subsection{Braille Cell Data}

We fuse DSBI~\cite{dsbi2019} (220 pages, 91,443 cells) and
Angelina~\cite{angel2021} (212 photos, 66,116 cells), plus 12 in-house
Gold pages, yielding \textbf{157,559} labelled cells.  DSBI annotates the
tight 2$\times$3 dot lattice while Angelina annotates the full cell; we
harmonise by expanding each DSBI box by $\alpha{=}0.35$ times the local
dot pitch.  Cleaning drops duplicate annotations, out-of-page boxes,
markout codes, and degenerate aspect ratios.  Page-group splits prevent
leakage (train 120,809 / val 21,095 / test 15,655).  Class imbalance is
handled via inverse-frequency loss weights, a per-class cap of 4,000, and
domain-balanced sampling.

\subsection{Fingertip Detection Data}

The fingertip detector is trained on three sources remapped to a single
\texttt{fingertip} class: (i)~\textbf{TI1K}~\cite{alam2019ti1k} (1,000 images;
thumb/index tip keypoints converted to YOLO boxes),
(ii)~\textbf{Roboflow Finger Tip Detection} (Universe; $\sim$538 images;
multi-finger tip boxes remapped to class~0), and (iii)~our
\textbf{Braille Fingertip Dataset} (60 LabelMe-annotated photos of fingers
on embossed Braille pages; split 48/6/6).  TI1K and Roboflow provide scale
and pose diversity; the Braille set adapts the detector to the target
reading domain.

\subsection{Braille Cell CNN}

A lightweight CNN classifies $64\times64$ greyscale cell crops into 64
Braille codes (three conv blocks + two FC layers).  Training uses
\texttt{AdamW}, cosine annealing, and label smoothing
($\varepsilon{=}0.05$).  The model reaches \textbf{94.7\%} top-1 on the
held-out test split and is exported to ONNX.

\subsection{Fingertip Detector}

We fine-tune YOLO v2.6 nano (\textit{yolo26n})~\cite{yolov8} on the
combined fingertip corpus.  Losses: BCE ($\lambda_\text{cls}{=}0.5$) for
confidence, CIoU + L1 ($\lambda_\text{box}{=}7.5$) for boxes (DFL omitted
for tight square tips).  Training runs 80 epochs at $640\times640$ with
ProgLoss warm-up and STAL one-to-one matching.  The model is UINT8-quantised
to ONNX ($\sim$3.5\,MB) for on-device inference, with FP32 as fallback.

\subsection{Cell Mapping and On-Device App}

A prescan builds an in-memory \textit{CellMap} (cell boxes, codes, Sinhala
characters).  Live frames are registered to the prescan via ORB+RANSAC
homography; the TipEMA-smoothed fingertip is hit-tested against expanded
cell boxes, and a 500\,ms dwell filter triggers character output.  A
Flutter app runs both ONNX models fully on-device (Android/iOS).

% ======================================================================
\section{Discussion}
% ======================================================================

\paragraph{Strengths.} The box-harmonisation step is a low-cost, annotation-free
solution to a real multi-source domain gap; it eliminates a systematic crop
misalignment that would otherwise degrade CNN accuracy.  The on-device
deployment removes the single most common barrier to adoption in
low-connectivity classrooms.

\paragraph{Limitations.} The current system has not yet been evaluated in a
formal user study with blind or visually impaired participants.  The Gold
in-house pages lack manually verified labels, so training with them awaits
annotation completion.  The cell detector was trained for only 1 epoch on
CPU for this submission (mAP50 0.125); full 80-epoch GPU training is
pending.

\paragraph{Future work.} We plan to (1)~complete GPU training and
report full detector mAP, (2)~add a text-to-speech engine for full Sinhala
narration, (3)~conduct a usability study with Grade 1 Braille students at a
Sri Lankan school for the blind, and (4)~extend support to Tamil Braille.

% ======================================================================
\section{Conclusion}
% ======================================================================

We have presented BrailleLens, an end-to-end on-device system for
interactive Braille learning with Sinhala output.  The pipeline integrates
multi-source Braille data through a principled harmonisation procedure,
trains a lightweight CNN and a domain-adapted fingertip detector, and
delivers real-time cell-level feedback via a cross-platform Flutter app.
The system runs entirely on consumer smartphones, making it practical for
deployment in resource-constrained educational settings.

% ======================================================================
\begin{ack}
The authors thank the University of Moratuwa, Department of Computer
Science and Engineering, for laboratory access and computing resources.
This work was carried out as part of the CS3501 Data Science and
Engineering Project course.
\end{ack}

% ======================================================================
\section*{References}
\begin{small}

[1] World Health Organization (2023). \textit{Blindness and vision impairment}.
\url{https://www.who.int/news-room/fact-sheets/detail/blindness-and-visual-impairment}.

[2] Antonacopoulos, A.\ \& Bridson, D.\ (1998). A robust braille recognition system.
\textit{Document Analysis Systems II}, pp.\ 533–545.

[3] Ovodov, I.G.\ (2021). Optical Braille recognition using object detection neural networks.
\textit{arXiv:2106.10765}.

[4] Nagata, A.\ et al.\ (2021). Angelina Braille reader: A smartphone-based Braille
reading aid. \textit{Proc.\ ICCHP}.

[5] Zhang, F.\ et al.\ (2020). MediaPipe Hands: On-device real-time hand tracking.
\textit{arXiv:2006.10214}.

[6] Redmon, J.\ \& Farhadi, A.\ (2018). YOLO v3: An incremental improvement.
\textit{arXiv:1804.02767}.

[7] Ultralytics (2024). YOLO v8 / YOLO v2.6 documentation.
\url{https://docs.ultralytics.com}.

[8] DSBI (2019). The Digitized Sinhala Braille Image dataset.
Department of Computer Science and Engineering, University of Moratuwa.

[9] Alam, M.M.\ et al.\ (2019). TI1K: A thousand-image fingertip dataset
with hand and tip annotations.
\url{https://github.com/MahmudulAlam/TI1K-Dataset}.

[10] Roboflow Universe. Finger Tip Detection dataset.
\url{https://universe.roboflow.com/first-leo0f/finger-tip-detection-i4xqf}.

\end{small}

%%%%%%%%%%%%%%%%%%%%%%%%%%%%%%%%%%%%%%%%%%%%%%%%%%%%%%%%%%%%
% NeurIPS checklist (same folder — no path prefix needed)
\newpage
\input{checklist}
%%%%%%%%%%%%%%%%%%%%%%%%%%%%%%%%%%%%%%%%%%%%%%%%%%%%%%%%%%%%

\end{document}
